\subsection{Phase-Based Constraint Resolution}
\label{sec:constraint_resolution}

\paragraph{Task Description.}
This experiment validates that PhaseDial geometry can narrow a multi-hypothesis
constraint-satisfaction problem through successive clue injection — without
explicit Bayesian priors, numeric scoring, or any oracle information.

Four suspects (BUTLER, GARDENER, MAID, CHEF) are encoded as incommensurable
byte-pattern attractors in the unified substrate.  Each suspect's samples
follow a distinct linear modular stride (3, 7, 11, 13 respectively), producing
maximally separated regions of the value/PhaseDial attractor landscape.

After ingestion, four stages of $4$ clue prompts are presented.
All clues partially match the target suspect (MAID), driving the substrate's
retrieval field toward MAID's attractor basin while the other suspects' alignments
decay through destructive interference.

\paragraph{Results.}
Figure~\ref{fig:constraint_resolution} shows four temporal snapshots of the
constraint-resolution process on the PhaseDial polar grid.  Each panel plots
suspects and the evidence centroid as dots at $(\theta_j, A_j)$, where
$\theta_j$ is phase angle (suspect assignment) and $A_j$ is retrieval alignment.

Mean final alignment with target (MAID): 0.000.
Mean final alignment averaged across all suspects: 0.000.
Target isolation ratio: 0.000.

\paragraph{Assessment.}
Partial isolation was achieved.  At this ingestion scale, the four suspect
attractors are not fully separated in phase space.  Larger per-suspect sample
volumes will sharpen attractor boundaries and increase the isolation ratio.

